\section{The Particular Judgment Immediately After Death}

Long before the heavens roll away and the graves give up their dead, every soul keeps a solitary appointment with Christ. The Epistle to the Hebrews states the law of mortal existence in twelve spare words: ``And as it is appointed unto men once to die, and after this the judgment'' (Heb 9:27). Between death and reckoning there is no second earthly probation, no reincarnation, and no dreamless sleep postponing the soul's destiny until the end of time. Ecclesiasticus already warned under the Old Covenant: ``it is easy before God in the day of death to reward every one according to his ways \dots\ and in the end of a man is the disclosing of his works'' (Ecclus 11:28--29).

The Gospels place immediate retribution directly before our eyes. In St.\ Luke's account of the rich man and Lazarus, neither soul waits in suspense for the end of the world: ``the beggar died and was carried by the angels into Abraham's bosom. And the rich man also died: and he was buried in hell'' (Luke 16:22). While the rich man's five brothers still live on earth with Moses and the prophets before them, Lazarus is already comforted, the rich man is already tormented in the flame, and between them ``there is fixed a great chaos: so that they who would pass from hence to you cannot, nor from thence come hither'' (Luke 16:25--28). Even more tender is our Lord's word from the Cross to the penitent thief: ``Amen I say to thee: This day thou shalt be with me in paradise'' (Luke 23:43). St.\ Paul speaks from the same certainty, desiring ``to be dissolved and to be with Christ, a thing by far the better'' (Phil 1:23) and longing ``to be absent rather from the body and to be present with the Lord'' (2 Cor 5:8). As St.\ Thomas Aquinas demonstrates against those who deferred heaven or hell until the last day, to be present with the Lord in contrast to walking by faith is to walk by sight---``to see God in His Essence, wherein consists eternal life.''\endnote{St.\ Thomas Aquinas, \textit{Summa Theologiae} III, q.\ 59, a.\ 5, ad 1, trans.\ Fathers of the English Dominican Province.}

Three major witnesses of the Church's magisterium seal this immediate destiny as Catholic dogma. In the Profession of Faith drafted by Pope Clement IV in 1267 and sworn in the name of Emperor Michael VIII Palaeologus at the Second Council of Lyons in 1274, the Church confessed that souls departing in charity without having completed worthy fruits of penance undergo purifying punishments (\textit{poenis purgatoriis}) after death; that souls free of every stain, either at once or after purification, are received immediately into heaven; and that souls departing in actual mortal sin, or in original sin alone, descend immediately into hell, though to unequal punishments according to their state.\endnote{Profession of Faith of Pope Clement IV (1267), sworn in the name of Emperor Michael VIII Palaeologus at the Second Council of Lyons (1274), DS 856--858.}

Sixty-two years later, after Pope John XXII had proposed as a private theological opinion---and retracted before his death---the view that separated souls of the just await the general resurrection before beholding the divine essence, Pope Benedict XII settled the question permanently in the Apostolic Constitution \textit{Benedictus Deus} (29 January 1336). He solemnly defined that, by God's general ordinance, the souls of all the saints who departed before Christ's Passion and of all the baptized faithful in whom nothing remained to be purged at death---or who have completed their purification afterward---enter heaven even before the resumption of their bodies and the general judgment. There they behold the divine essence directly, clearly, and face to face by an intuitive vision, with no creature mediating as the object seen, and in that vision they enjoy true blessedness and eternal rest.\endnote{Pope Benedict XII, Apostolic Constitution \textit{Benedictus Deus} (29 January 1336), DS 1000--1001.}

In the same constitution, Benedict XII solemnly defined that souls departing in actual mortal sin descend immediately after death into hell to suffer infernal torments, while reaffirming that all men will nonetheless appear with their bodies before Christ's tribunal on the Last Day to render an account of their own deeds (DS 1002). A century later, the Ecumenical Council of Florence in the Decree of Union \textit{Laetentur caeli} (6 July 1439) reaffirmed this exact threefold immediate destiny in common between East and West, defining further that the blessed behold the Triune God clearly as He is, yet one more perfectly than another according to the diversity of their merits.\endnote{Council of Florence (1439), Decree of Union \textit{Laetentur caeli} (DS 1304--1306).}

This solemnly defined immediate retribution presupposes the judicial act by which Christ assigns each soul its state: the \textit{Particular Judgment}, taught steadily by the Ordinary and Universal Magisterium. The \textit{Roman Catechism} instructs pastors to place it first before the faithful:

\begin{quote}
The first, when each one departs this life; for he is instantly placed before the judgment seat of God, where all that he had ever done, or spoken, or thought, during life, shall be subjected to the most rigid scrutiny; and this is called the particular judgment.\endnote{\textit{Catechism of the Council of Trent}, Part I, Creed Article VII, trans.\ Donovan, p.\ 62.}
\end{quote}

The \textit{Catechism of the Catholic Church} summarizes this same faith: death closes earthly life as the time for accepting or refusing the grace manifested in Christ, and at the very instant of death each person receives eternal retribution in his immortal soul through a particular judgment referring his life directly to Christ---leading either into heaven's beatitude, immediately or through purification, or into immediate and everlasting damnation.\endnote{\textit{Catechism of the Catholic Church}, \S\S 1021--1022.}

And the \textit{Baltimore Catechism No.\ 3} adds three plain points for daily remembrance: the particular judgment takes place ``in the place where each person dies, and the soul will go immediately to its reward or punishment''; its sentence ``will not be changed at the general judgment, but it will be repeated and made public to all''; and its three outcomes are ``Heaven, Purgatory, and Hell.''\endnote{\textit{A Catechism of Christian Doctrine, No.\ 3} (Baltimore Catechism No.\ 3; New York: Benziger Brothers, 1921), Qq.\ 1372--1378, pp.\ 306--307.}

Theologically, this encounter is neither a protracted earthly trial nor a second chance to reverse a life's choice. As St.\ Thomas teaches, while united to its mortal body the human will remains flexible; once death separates soul from body, the soul enters an unchangeable state fixed in the charity or refusal with which it died.\endnote{Aquinas, \textit{ST} III, q.\ 59, a.\ 5, corpus and ad 1; cf.\ \textit{ST} III, q.\ 86, a.\ 1, corpus.} In one instantaneous flash of Christ's truth, every earthly disguise collapses. Excuses, self-flattery, minimized sins, and forgotten omissions vanish; the soul sees its thoughts, words, deeds, graces, and final orientation exactly as God sees them, and its own illuminated conscience ratifies Christ's sentence.

Three destinies alone open from that instant: immediate admission to the beatific vision for the soul departing in grace free of all stain and temporal debt; entrance into Purgatory for the soul departing in sanctifying grace yet still bearing venial faults, disordered attachments, or unpaid temporal satisfaction; or immediate, irrevocable descent into hell for the soul departing in unrepented actual mortal sin.

The saints turned this certainty into a daily school of vigilance. In Consideration XXIV of \textit{Preparation for Death}, St.\ Alphonsus Maria de Liguori notes that it is the common opinion of theologians that Christ Himself judges the soul at the instant of death in the very place where it leaves the body, laying open two books: ``There shall be two books, the Gospel and conscience. In the Gospel shall be read what the accused should have done, and in his conscience, what he has done.''\endnote{St.\ Alphonsus Maria de Liguori, \textit{Preparation for Death}, Consideration XXIV, First and Second Points, trans.\ by a Catholic clergyman (Dublin: James Duffy, 1843), pp.\ 289, 293.} Rebuking ``foolish security in so great a peril'' (\textit{Quae est ista stulta securitas in discrimine tanto?}, from St.\ Thomas of Villanova), St.\ Alphonsus urges St.\ Bernard's watchword:

\begin{quote}
At present it depends on us to choose whatever sentence we please. \dots\ We must adjust our accounts before judgment. ``Before judgment, prepare thee justice''---Eccl., xviii.\ 19. \dots\ Let us then say with St.\ Bernard: ``\textit{Volo judicatus praesentari, non judicandus}.'' O my Judge, I wish to be judged and punished during life, which is a time of mercy and pardon; for, after death shall be the time of justice.\endnote{St.\ Alphonsus Liguori, \textit{Preparation for Death}, Consideration XXIV, Third Point, pp.\ 298--299.}
\end{quote}

St.\ Francis de Sales leads the soul to the exact same resolution in his Sixth Meditation on Judgment:

\begin{quote}
Ah! I will judge myself now that I may not be judged then! I will examine my conscience, and condemn myself; I will accuse myself, and amend my life, that the eternal Judge may not condemn me on that dreadful day. I will, therefore, confess my sins, and receive all necessary advice.\endnote{St.\ Francis de Sales, \textit{Introduction to a Devout Life} I.14 (Sixth Meditation: On Judgment), New York: Fr.\ Pustet, p.\ 37.}
\end{quote}

Every sincere nightly examination, every humble confession, and every act of restitution anticipates that solitary instant while the door of mercy stands wide open.
